Atmospheric refraction raises the Sun above its horizon, at its limit, at
sunrise and sunset, approximately by an apparent disk, as shown in the figure.
\ioaafig{2022_TH_Q02-s1.pdf}
The upper arc of the real solar disk rises after the moment when the whole
apparent disk is already visible above the horizon. The official explanation
then states that sunrise and sunset are seen later than they actually occur,
making the day a little longer.
% sic: The source says sunrise is seen later; the preceding sentence describes an earlier apparent sunrise.
Under these conditions,
\begin{align*}
\rho_{r,\max}&=\theta_S,\\
\sin\frac{\theta_S}{2}&=\frac{R_S}{d}\approx\frac{\theta_S}{2},\\
\theta_S&=\frac{2R_S}{d}=\frac{2\times6.96\times10^5}{15\times10^7}
=0.0098\,\mathrm{rad}.
\end{align*}
% sic: The printed decimal 0.0098 is inconsistent with the preceding ratio and the following value in arcminutes.
Using $1\,\mathrm{rad}=180\times60'/3.14$, the source gives
\[
\theta_S=0.0098\frac{180\times60'}{3.14}\approx33.7'=\rho_{r,\max}.
\]
Thus the official answer is (d).
